\section{The Dock}

One cell of the mesh is a dock, a here-and-now, the smallest place where anything happens. Its shape is the 24-cell, with twenty-four corners and twenty-four faces, equal to its own dual. Each dock meets exactly twenty-four neighbors, one across each face.

Out of each dock run twenty-four directions, one toward each neighbor, and each direction is a site. A site is the outward normal through one of the dock's twenty-four faces, pointing straight to the dock waiting on the other side. Stand at the center of a dock and look out, and the sites are the twenty-four evenly spread directions you could step along, as symmetric a spray of arrows as four-dimensional space allows. The twenty-four sites pair off into twelve lines, each a direction together with its opposite, and the line, not the single site, is what the law acts on.

A single tone lives at each site, so a dock holds twenty-four tones at once, one per site. Its full state is a point in a space of three to the twenty-fourth, about two hundred eighty billion possibilities, and by the structure of the directions that state splits into three octets, a sensory part and two of opposite handedness, woven into one rich quality. A single dock already carries far more structure than any single number.

A free tone, one not caught in a collision, has its value handed one dock per beat along its site, never spreading out. That straight passing is momentum, and a tone's momentum is simply which site it occupies. Step back and the dock is two things at once. It is the single tile from which the whole mesh is laid, the unit cell of the world repeated without end. And it is the smallest place where the dynamics happens, one here-and-now where tones meet their neighbors across the twenty-four faces and are passed on. Everything the rest of the paper builds, a particle, a force, a bound knot of matter, is made of these and only these, vast numbers of them woven together and stepped in time.
